% Part: history
% Chapter: biographies
% Section: alfred-tarski

\documentclass[../../../include/open-logic-section]{subfiles}

\begin{document}

\olfileid{his}{bio}{tar}

\olsection{Alfred Tarski}

\olphoto{tarski-alfred}{Tarski}

阿尔弗雷德·Tarski于1901年1月14日出生于波兰华沙（当时属俄罗斯帝国）。Tarski被形容为“拿破仑式的人物”，他精力旺盛、健谈且充满激情。他的这种活力常常体现在课堂上——有一次讲课时扔弃烟头，他引燃了废纸篓，因而被禁止再在那栋楼里讲课。

Tarski从小就渴求知识。尽管他后来告诉学生，自己之所以学逻辑是因为那是他唯一得过 B 的课程，但他的高中成绩单却显示他门门得 A——连逻辑课也不例外。1918年至1924年，他就读于华沙大学。Tarski最初打算修读生物学，但由于当时的华沙大学正是华沙逻辑与哲学学派的中心，他对数学、哲学和逻辑产生了兴趣。1924年，Tarski在 Stanisław Leśniewski（莱什涅夫斯基）的指导下获得博士学位。

在1939年移居美国之前，Tarski在华沙担任中学教师期间完成了他的一些最重要的工作。他关于逻辑后承与逻辑真理的论著便是这段时间写就的。1939年，Tarski正在美国进行巡回演讲。在他访美期间，德国入侵波兰；由于犹太血统，Tarski无法返回。他的妻儿在战争结束前一直留在波兰，战后才得以移居美国与他团聚。Tarski先后在哈佛大学、纽约市立学院和普林斯顿高等研究院任教，最后任职于加州大学伯克利分校。在那里，他创立了“逻辑与科学方法论”的多学科项目。Tarski于1983年10月26日去世，享年82岁。

\begin{reading}
如需了解Tarski生平的更多细节，可参阅传记 \emph{Alfred Tarski: Life and Logic} \citep{Feferman2004}。Tarski关于逻辑后承与逻辑真理的开创性论著的英译本见 \citep{Tarski1983}。Tarski的所有原著已汇编为四卷本系列 \citep{Tarski1981}。
\end{reading}

\end{document}
